\usemodule[memos]
\usemodule[basicexam]
\usemodule[setups]

\enablemode[no-setup-main]
\loadsetups[i-basicexam.xml]

\setupinteraction[state=start]
\setuptextrules[style=\tf\ss,before=,after=,rulethickness=.5pt]

\setuptyping[style=\ttxx\setuplocalinterlinespace[line=2ex],
             option={context},
             before={\vbox{\hrule height .5pt\par
                     \textrule{代码示例}%
                     \hrule height .5pt\vrule width0pt depth 4pt}\par},
              after={\vbox{\hrule height .5pt\strut}},
             ]

\starttext

\title{basicexam 模块使用说明书}

\blank[big]

\placecontent

\chapter{模块概述}

\section{简介}

basicexam 是一个为 ConTeXt LMTX 设计的试卷制作模块，提供了完整的试卷排版解决方案。
该模块支持多种题型，包括选择题、填空题、简答题、完形填空等，并具有答案管理、分数统计等功能。

\section{主要特性}

\startitemize[packed]
\item 多种题型支持：选择题、填空题、简答题、完形填空等
\item 教师/学生模式切换，控制答案显示
\item 无限层级嵌套题目支持
\item 自动分数统计与答案收集
\item 灵活的样式定制选项
\item 内置多种符号标记
\item 作文格与写作区域
\item 材料阅读与引用系统
\stopitemize

\section{系统要求}

该模块需要 {\bf ConTeXt LMTX (LuaMetaTeX)} 环境，不支持 MkIV 版本。
如果在不兼容的环境下运行，模块会提示错误并终止编译。

\chapter{模块加载与参数}

\section{加载模块}

\starttyping
\usemodule[basicexam]
\stoptyping

\section{模块参数}

模块加载时可设置以下参数：

\starttabulate[|l|l|p|]
\HL
\NC 参数 \NC 可选值 \NC 说明 \NC\NR
\HL
\NC mode \NC \NC 运行模式：\NC\NR
\NC      \NC student  \NC 学生模式，隐藏答案 \NC\NR
\NC      \NC teacher  \NC 教师模式，显示答案 \NC\NR
\NC      \NC check    \NC 检查模式，显示答案并输出调试信息 \NC\NR
\HL
\NC layout \NC \NC 布局设置：\NC\NR
\NC        \NC twoup   \NC 双页排版（每页两版）\NC\NR
\NC        \NC none    \NC 默认单页布局 \NC\NR
\HL
\NC footer \NC \NC 页脚设置：\NC\NR
\NC        \NC default \NC 显示科目和页码 \NC\NR
\NC        \NC none    \NC 无页脚 \NC\NR
\HL
\NC libpath \NC 路径 \NC SQLite 库文件路径（用于数据存储）\NC\NR
\HL
\NC showanswer \NC showanswer \NC 显示答案 \NC\NR
\NC showpoint  \NC showpoint  \NC 显示分数 \NC\NR
\NC showscript \NC showscript \NC 显示答题区域 \NC\NR
\HL
\stoptabulate

\subsection{示例}

\starttyping
% 教师模式，显示答案
\usemodule[basicexam][mode=teacher]

% 学生模式，隐藏答案
\usemodule[basicexam][mode=student]

% 检查模式，调试输出
\usemodule[basicexam][mode={teacher,check}]

% 双页排版
\usemodule[basicexam][layout=twoup]
\stoptyping

\section{模式控制}

\subsection{全局模式}

通过模块参数设置全局模式：

\starttyping
\usemodule[basicexam][mode=teacher]  % 教师模式
\usemodule[basicexam][mode=student]  % 学生模式
\stoptyping

\subsection{局部控制}

在文档中动态切换：

\starttyping
\showanswer   % 显示答案
\hideanswer   % 隐藏答案
\showpoint    % 显示分数
\hidepoint    % 隐藏分数
\stoptyping

\subsection{单个题目控制}

在题目级别控制显示：

\starttyping
\startquestion[showanswer=true,showpoint=true]
  这道题会显示答案和分数。
\stopquestion

\startquestion[showanswer=false,showpoint=false]
  这道题会隐藏答案和分数。
\stopquestion
\stoptyping

\section{样式定制}

\subsection{全局设置}

使用 \type{\setupquestion} 进行全局设置：

\startbuffer
\setupquestion[
  answer={默认答案},
  point=1,
  showanswer=true,
  showpoint=true,
  style=\ss,
  color=black,
]
\stopbuffer

\typebuffer

\subsection{选择题样式}

\startbuffer
\setupchoice[
  answerstyle=\ss,
  answercolor=red,
  numberconversion=A,
  option={A,packed,joinedup,nowhite},
]
\stopbuffer

\typebuffer

\subsection{答案样式}

\startbuffer
\setupanswer[
  style=\ss\it,
  color=darkblue,
  answerlabel={{答案}},
  pointlabel={，分},
  before=\blank[halfline]\startnarrower,
  after=\blank[halfline]\stopnarrower,
]
\stopbuffer

\typebuffer

\chapter{试卷标题}

\section{papertitle 命令}

试卷标题是试卷的重要组成部分，通常包含考试名称、科目、年级、考试时间等信息。
模块提供 \type{\definepapertitle} 和 \type{\makepapertitle} 命令来创建标准化的试卷标题。

\subsection{基本用法}

首先使用 \type{\definepapertitle} 定义标题样式，然后使用 \type{\setuppapertitle} 设置内容，最后用 \type{\makepapertitle} 输出：

\startbuffer
\definepapertitle[myexam][list={secret,title,subject,information}]
\setuppapertitle[myexam][
  secret={绝密 ★ 启用前},
  title={2024 年普通高等学校招生全国统一考试},
  subject={数学},
  information={总分:150 分，考试时间:120 分钟}]
\makepapertitle[myexam]
\stopbuffer

\typebuffer

\getbuffer

\subsection{标题元素}

\type{list} 参数控制标题中显示哪些元素及其顺序。默认列表为：\type{secret,title,subject,information,signinformation,notice}。

\starttabulate[|l|p|]
\HL
\NC 元素 \NC 说明 \NC\NR
\HL
\NC secret \NC 密级标识（如"绝密 ★ 启用前"） \NC\NR
\NC title \NC 主标题 \NC\NR
\NC subject \NC 科目名称（自动添加 subjectsuffix 后缀） \NC\NR
\NC information \NC 考试信息（如总分、时间） \NC\NR
\NC signinformation \NC 签名信息区域 \NC\NR
\NC notice \NC 注意事项 \NC\NR
\HL
\stoptabulate

每个元素都可以单独设置样式、颜色和对齐方式，例如：
\type{secretstyle}、\type{secretcolor}、\type{secretalign} 等。

\subsection{其他参数}

\starttabulate[|l|l|p|]
\HL
\NC 参数 \NC 默认值 \NC 说明 \NC\NR
\HL
\NC alternative \NC default \NC 标题模板 \NC\NR
\NC list \NC \NC 显示元素列表 \NC\NR
\NC subjectsuffix \NC 学科试卷 \NC 科目后缀 \NC\NR
\NC style \NC \NC 全局字体样式 \NC\NR
\NC color \NC \NC 全局字体颜色 \NC\NR
\NC align \NC center \NC 全局对齐方式 \NC\NR
\HL
\stoptabulate

\section{自定义标题样式}

可以通过继承已有样式创建新的标题样式：

\startbuffer
\definepapertitle[mytitle][myexam]
\setuppapertitle[mytitle][title={我的试卷},style=\ssa\bf]
\makepapertitle[mytitle]
\stopbuffer

\typebuffer

\start\getbuffer\stop

\chapter{选择题}

\section{基本用法}

选择题有两种语法格式：{\bf 环境格式}和{\bf 快捷格式}。

\subsection{环境格式}

使用 \type{\startquestion}...\type{\stopquestion} 和 \type{\startchoice}...\type{\stopchoice} 环境：

\startbuffer
\startquestion
  这是选择题的题干内容。
  \startchoice
    \startcitem 选项一 \stopcitem
    \startcitem 选项二 \stopcitem
    \startcitem 选项三 \stopcitem
    \startcitem 选项四 \stopcitem
  \stopchoice
\stopquestion
\stopbuffer

\typebuffer \getbuffer

\subsection{快捷格式}

使用 \type{\question} 和 \type{\choice} 命令：

\startbuffer
\question{这是选择题的题干内容。
  \choice{选项一,选项二,选项三,选项四}
}
\stopbuffer

\typebuffer \getbuffer

\section{参数说明}

\showsetup[setupchoice]

\starttabulate[|l|l|p|]
\HL
\NC 参数 \NC 默认值 \NC 说明 \NC\NR
\HL
\NC inlinechoice \NC false \NC 是否为内联选择题 \NC\NR
\NC inlineleft \NC \NC 内联选项左侧内容 \NC\NR
\NC inlineright \NC \NC 内联选项右侧内容 \NC\NR
\NC correctsymbol \NC \NC 正确答案标记符号 \NC\NR
\NC separator \NC \NC 选项分隔符 \NC\NR
\NC textstyle \NC \NC 选项文本样式 \NC\NR
\NC textcolor \NC \NC 选项文本颜色 \NC\NR
\NC showanswer \NC \NC 是否显示答案 \NC\NR
\NC showpoint \NC \NC 是否显示分数 \NC\NR
\NC answer \NC \NC 答案内容 \NC\NR
\NC answerstyle \NC \NC 答案样式 \NC\NR
\NC answercolor \NC \NC 答案颜色 \NC\NR
\NC answerlabel \NC \NC 答案标签 \NC\NR
\NC point \NC \NC 题目分数 \NC\NR
\NC pointstyle \NC \NC 分数样式 \NC\NR
\NC pointcolor \NC \NC 分数颜色 \NC\NR
\NC pointlabel \NC \NC 分数标签 \NC\NR
\HL
\stoptabulate

\section{进阶用法}

\subsection{标记正确答案}

使用 \type{[*]} 标记正确答案：

\startbuffer
\startquestion[showanswer=true]
  这是选择题的题干内容。
  \startchoice
    \startcitem    选项一 \stopcitem
    \startcitem    选项二 \stopcitem
    \startcitem[*] 选项三 \stopcitem
    \startcitem    选项四 \stopcitem
  \stopchoice
\stopquestion

\question[showanswer=true]{这是选择题的题干内容。
  \choice{选项一,选项二,{[*]选项三},选项四}
}
\stopbuffer

\typebuffer \getbuffer

\subsection{多选题}

支持标记多个正确答案：

\startbuffer
\startquestion[showanswer=true]
  请选择所有正确的选项。
  \startchoice
    \startcitem[*] 选项一 \stopcitem
    \startcitem    选项二 \stopcitem
    \startcitem[*] 选项三 \stopcitem
    \startcitem    选项四 \stopcitem
  \stopchoice
\stopquestion
\stopbuffer

\typebuffer \getbuffer

\subsection{选项列数控制}

通过可选参数控制选项的显示列数：

\startbuffer
\startquestion
  选项以8列显示。
  \startchoice[8]
    \startcitem A \stopcitem
    \startcitem B \stopcitem
    \startcitem C \stopcitem
    \startcitem D \stopcitem
    \startcitem E \stopcitem
    \startcitem F \stopcitem
    \startcitem G \stopcitem
    \startcitem H \stopcitem
  \stopchoice
\stopquestion

\question{选项以8列显示。
  \choice[8]{A,B,C,D,E,F,G,H}
}
\stopbuffer

\typebuffer \getbuffer

\subsection{内联选择题}

内联选择题将选项显示在题干行内，适合简短的选择题：

\startbuffer
\startquestion
  这是内联选择题。\relax
  \startinlinechoice
    \startcitem 选项一 \stopcitem
    \startcitem 选项二 \stopcitem
    \startcitem 选项三 \stopcitem
    \startcitem 选项四 \stopcitem
  \stopinlinechoice
\stopquestion

\question{这是内联选择题。\relax
  \inlinechoice{选项一,选项二,选项三,选项四}
}
\stopbuffer

\typebuffer \getbuffer

\section{自定义样式}

使用 \type{\definechoice} 创建自定义选择题样式：

\startbuffer
\definechoice[mychoice][
  option={text,6,packed,joinedup,nowhite},
  inlinechoice=true,
  correctsymbol={\symbol{markcirclecheck}},
  inlineleft={【\nobreak},
  inlineright={\nobreak 】},
  answerstyle=\ss,
  separator={\removeunwantedspaces,\space},
]

\startquestion[showanswer=true]
  使用自定义选择题样式。
  \startmychoice
    \startcitem[*] 选项一 \stopcitem
    \startcitem    选项二 \stopcitem
    \startcitem    选项三 \stopcitem
    \startcitem    选项四 \stopcitem
  \stopmychoice
\stopquestion
\stopbuffer

\typebuffer \getbuffer

\chapter{题目与问题}

\section{question 环境}

\type{question} 环境用于创建独立题目：

\startbuffer
\startquestion[option={n,packed,joinedup,nowhite}]
这是题目的主要内容。
\stopquestion

\question{这是题目的主要内容。}
\stopbuffer

\typebuffer \getbuffer

\section{problem 环境}

\type{problem} 环境用于创建子问题列表：

\startbuffer
\startquestion
这是大题的题干。
  \startproblem
  \startpitem 这是第一个小题。\stoppitem
  \startpitem 这是第二个小题。\stoppitem
  \stopproblem
\stopquestion

\question{这是大题的题干。
  \problem[textstyle=ss]
    {\pitem{这是第一个小题。}
     \pitem{这是第二个小题。}
  }
}
\stopbuffer

\typebuffer \getbuffer

\section{题目嵌套}

支持无限层级的题目嵌套：

\startbuffer
\startquestion
这是第一层题目。
  \startquestion
  这是第二层题目。
    \startquestion
    这是第三层题目。
    \stopquestion
  \stopquestion
\stopquestion
\stopbuffer

\typebuffer \getbuffer

\section{question 参数}

\showsetup[setupquestion]

\chapter{答案与解析}

\section{基本用法}

\type{answer} 环境用于添加答案和解析：

\startbuffer
\startquestion[point=4,showpoint=true,answer=A]
这是选择题的题干。
  \startchoice
    \startcitem[*] 选项A \stopcitem
    \startcitem    选项B \stopcitem
    \startcitem    选项C \stopcitem
    \startcitem    选项D \stopcitem
  \stopchoice
  \startanswer
    这是对该题目的详细解析说明。
    答案选择A的原因是……
  \stopanswer
\stopquestion
\stopbuffer

\typebuffer \getbuffer

\section{参数说明}

\starttabulate[|l|l|p|]
\HL
\NC 参数 \NC 默认值 \NC 说明 \NC\NR
\HL
\NC showanswer \NC \NC 是否显示答案 \NC\NR
\NC showpoint \NC \NC 是否显示分数 \NC\NR
\NC showscript \NC \NC 是否显示答题区域 \NC\NR
\NC scriptcontent \NC \NC 答题区域提示内容 \NC\NR
\NC style \NC \ss\it \NC 答案字体样式 \NC\NR
\NC color \NC \NC 答案颜色 \NC\NR
\NC answerlabel \NC 答案 \NC 答案标签 \NC\NR
\NC pointlabel \NC ,分 \NC 分数标签 \NC\NR
\HL
\stoptabulate

\section{进阶用法}

\subsection{答题区域}

使用 \type{showscript=true} 添加答题书写区域：

\startbuffer
\startquestion[point=4,showpoint=true]
这是需要书写答案的题目。
\startanswer[showscript=true,frame=off,
             scriptcontent={\centerline{\lightgray\ss 请在此处书写答案}\par
                            \null\vskip3\baselineskip\relax}]
\stopanswer
\stopquestion
\stopbuffer

\typebuffer \getbuffer

\subsection{分数设置}

在题目中设置分数：

\startbuffer
\question[point=5,showpoint=true]{这道题5分。\choice{A,B,C,D}}

\startquestion[point=10,showpoint=true]
这道大题共10分。
\startproblem
  \startpitem[point=3,showpoint=true] 第一小题3分。\stoppitem
  \startpitem[point=4,showpoint=true] 第二小题4分。\stoppitem
  \startpitem[point=3,showpoint=true] 第三小题3分。\stoppitem
\stopproblem
\stopquestion
\stopbuffer

\typebuffer \getbuffer

\subsection{score 命令}

\type{\score} 命令用于在答案中标注各要点的得分，方便评分者了解分值分布。
该命令通常与 \type{\resetscore} 配合使用，后者用于重置累计分数。

在答案中使用 \type{\score{分数}} 标注得分点：

\startbuffer
\startanswer
  答案要点一 \score{3}
  答案要点二 \score{4}
  答案要点三 \score{3}
  \resetscore
  总分 \score[type=dotfills]{10}
\stopanswer
\stopbuffer

\typebuffer \getbuffer

score 命令参数：

\starttabulate[|l|l|p|]
\HL
\NC 参数 \NC 默认值 \NC 说明 \NC\NR
\HL
\NC type \NC \NC 显示类型：dotfills、normal \NC\NR
\NC style \NC \NC 数字样式 \NC\NR
\NC color \NC \NC 数字颜色 \NC\NR
\NC before \NC \NC 分数前内容 \NC\NR
\NC after \NC \NC 分数后内容 \NC\NR
\HL
\stoptabulate

\subsection{答案输出}

答案输出功能用于在试卷末尾或指定位置生成答案汇总表。
使用 \type{\typeanswer} 命令可以按不同格式输出答案。

\startbuffer
% 列表格式（默认）
\typeanswer[start=1,total=10]

% 表格格式
\typeanswer[alternative=table,start=1,total=30,rows=6,columns=5]

% 网格格式
\typeanswer[alternative=grid,start=1,total=10,rows=5,columns=2]

% 直接访问单个答案
\typeanswer[alternative=direct]{Q-5}
\stopbuffer

\typebuffer

typeanswer 命令参数：

\starttabulate[|l|l|p|]
\HL
\NC 参数 \NC 默认值 \NC 说明 \NC\NR
\HL
\NC alternative \NC list \NC 输出格式：list、table、grid、direct \NC\NR
\NC start \NC 1 \NC 起始题号 \NC\NR
\NC stop \NC \NC 结束题号 \NC\NR
\NC total \NC \NC 题目总数（替代 stop） \NC\NR
\NC chapter \NC \NC 章节编号 \NC\NR
\NC rows \NC \NC 表格行数 \NC\NR
\NC columns \NC \NC 表格列数 \NC\NR
\NC separator \NC \NC 分隔符（文本格式） \NC\NR
\NC left \NC \NC 答案左侧内容 \NC\NR
\NC right \NC \NC 答案右侧内容 \NC\NR
\HL
\stoptabulate

\chapter{填空题}

\section{基本用法}

\type{\fillin} 命令用于在题目中创建填空项：

\startbuffer
\startquestion[showanswer=true,point=2,showpoint=true]
  Lorem ipsum dolor sit amet,
  \fillin{填空答案}
  consectetur adipiscing elit.
\stopquestion
\stopbuffer

\typebuffer \getbuffer

\section{参数说明}

\showsetup[setupfillin]

\starttabulate[|l|l|p|]
\HL
\NC 参数 \NC 默认值 \NC 说明 \NC\NR
\HL
\NC style \NC \NC 填空内容样式 \NC\NR
\NC color \NC \NC 填空内容颜色 \NC\NR
\NC before \NC \NC 填空前内容 \NC\NR
\NC after \NC \NC 填空后内容 \NC\NR
\NC left \NC \NC 左侧边界 \NC\NR
\NC right \NC \NC 右侧边界 \NC\NR
\NC empty \NC yes \NC 空白显示方式：yes、number、none \NC\NR
\NC alternative \NC default \NC 填空样式：default、frame、simple \NC\NR
\NC n \NC \NC 下划线长度（字符数） \NC\NR
\NC width \NC \NC 填空宽度 \NC\NR
\HL
\stoptabulate

\section{进阶用法}

\subsection{填空样式}

通过 \type{empty} 参数控制填空显示方式：

\startbuffer
\setupquestion[showanswer=false,point=2,showpoint=false]

\startquestion[showanswer=false]
  不同填空样式示例：
  \startproblem
    {\setupfillin[empty=number,n=5]
     编号样式：Lorem ipsum \fillin{答案}}
    
    {\setupfillin[empty=yes,n=5]
     隐藏样式：Lorem ipsum \fillin{答案}}
    
    {\setupfillin[empty=none,n=5]
     显示样式：Lorem ipsum \fillin{答案}}
  \stopproblem
\stopquestion
\stopbuffer

\typebuffer \getbuffer

\subsection{判断题}

填空题也可用于判断题：

\startbuffer
\setupfillin[empty=none,style=\tta]
\startquestion
判断下列说法是否正确。
  \startproblem
  \startpitem \fillin{T} 这是正确的说法。\stoppitem
  \startpitem \fillin{F} 这是错误的说法。\stoppitem
  \startpitem 这是正确的说法。\hfill\fillin{\symbol{markcheck}}\stoppitem
  \startpitem 这是错误的说法。\hfill\fillin{\symbol{markcross}}\stoppitem
  \stopproblem
\stopquestion
\stopbuffer

\typebuffer \getbuffer

\chapter{材料阅读}

\section{基本用法}

材料阅读题是试卷中常见的题型，通常包含一篇阅读材料和若干相关问题。
\type{material} 环境用于放置阅读材料，支持设置标题、作者、来源等信息。

\startbuffer
\setupmaterial[title][style=\ssa]
\setupmaterial[author][style=\tta]
\setupmaterial[source][style=ss]

\startmaterial[title={材料标题},author={作者名},source={来源}]
  这是阅读材料的内容。
  可以包含多段文字。
  \indicator{需要提问的内容一}
  \indicator{需要提问的内容二}
\stopmaterial

\startquestion
  \indicator{需要提问的内容一}
  \choice{选项A,选项B,选项C,选项D}
\stopquestion

\startquestion
  \indicator{需要提问的内容二}
  \choice{选项A,选项B,选项C,选项D}
\stopquestion
\stopbuffer

\typebuffer \getbuffer

\section{参数说明}

\showsetup[setupmaterial]

\starttabulate[|l|l|p|]
\HL
\NC 参数 \NC 默认值 \NC 说明 \NC\NR
\HL
\NC title \NC \NC 材料标题 \NC\NR
\NC author \NC \NC 作者 \NC\NR
\NC source \NC \NC 来源 \NC\NR
\NC style \NC \NC 正文字体样式 \NC\NR
\NC color \NC \NC 正文字体颜色 \NC\NR
\NC before \NC \NC 材料前内容 \NC\NR
\NC after \NC \NC 材料后内容 \NC\NR
\HL
\stoptabulate

\section{进阶用法}

\subsection{indicator 命令}

\type{\indicator} 命令用于标记材料中的提问点，实现材料与题目的关联。
标记后的内容会在材料中以特殊样式显示，同时在题目中引用。

\startitemize[packed]
\item 在材料中使用 \type{\indicator{标记名}} 标记关键内容
\item 在题目中使用 \type{\indicator{标记名}} 引用该内容
\item 编译时自动建立关联关系
\stopitemize

\startbuffer
\startmaterial[title={示例材料}]
  文章内容，其中\indicator{重点内容一}需要提问，
  还有\indicator{重点内容二}也需要提问。
\stopmaterial
\stopbuffer

\typebuffer \getbuffer

\chapter{诗词排版}

\section{基本用法}

\type{verse} 环境用于排版诗词、歌词等韵文内容：

\startbuffer
\startverse[title=临江仙 $\cdot$ 滚滚长江东逝水,author=明 $\cdot$ 杨慎,sample={白髮漁樵江渚上 慣看秋月春風}]
滾滾長江東逝水浪花淘盡英雄\\
是非成敗轉頭空\\
青山依舊在 幾度夕陽紅\\
白髮漁樵江渚上 慣看秋月春風\\
一壺濁酒喜相逢\\
古今多少事 都付笑談中
\stopverse
\stopbuffer

\typebuffer \getbuffer

\section{参数说明}

\showsetup[setupverse]

\starttabulate[|l|l|p|]
\HL
\NC 参数 \NC 默认值 \NC 说明 \NC\NR
\HL
\NC title \NC \NC 诗词标题 \NC\NR
\NC author \NC \NC 作者信息 \NC\NR
\NC source \NC \NC 来源信息 \NC\NR
\NC sample \NC \NC 用于计算缩进的样本文本 \NC\NR
\NC skip \NC both \NC 缩进方式：left、right、both、none \NC\NR
\NC align \NC \NC 对齐方式 \NC\NR
\NC style \NC \NC 字体样式 \NC\NR
\NC color \NC \NC 文字颜色 \NC\NR
\NC before \NC \NC 环境前内容 \NC\NR
\NC after \NC \NC 环境后内容 \NC\NR
\HL
\stoptabulate

\section{进阶用法}

可以分别设置标题、作者、来源的样式：

\startbuffer
\setupverse[sample={滾滾長江東逝水浪花淘盡英雄},skip=left,before=\blank,after=\blank]
\setupverse[title][align=center,style=bf,skip=both]
\setupverse[author][align=flushright,style=it,skip=both]
\setupverse[source][align=flushright,style=it,skip=both]
\stopbuffer

\typebuffer

\chapter{完形填空}

\section{基本用法}

\type{close} 环境用于创建完形填空题，在文章中嵌入选项供学生选择。
每个选项使用 \type{\closechoice} 命令标记，正确答案用 \type{[*]} 标识。

\startbuffer
\startclose[showanswer=true,showpoint=true,point=1.5,answer=nil]
On Oct. 11, hundreds of runners competed in a cross-country race.
Melanie Bailey should have \closechoice[designed,followed,changed,finished] the
course earlier than she did. Her \closechoice[delay,chance,trouble,excuse] came
because she was carrying a \closechoice[judge,volunteer,classmate,competitor]
across the finish line.
\stopclose
\stopbuffer

\typebuffer \getbuffer

\section{参数说明}

\starttabulate[|l|l|p|]
\HL
\NC 参数 \NC 默认值 \NC 说明 \NC\NR
\HL
\NC reset \NC false \NC 是否重置题号 \NC\NR
\NC n \NC 8 \NC 下划线长度 \NC\NR
\NC width \NC 1.5em \NC 题号宽度 \NC\NR
\NC answerwidth \NC 4em \NC 答案区域宽度 \NC\NR
\NC point \NC \NC 每题分数 \NC\NR
\NC answer \NC \NC 默认答案 \NC\NR
\NC showanswer \NC \NC 是否显示答案 \NC\NR
\NC showpoint \NC \NC 是否显示分数 \NC\NR
\NC style \NC \NC 文本样式 \NC\NR
\NC color \NC \NC 文本颜色 \NC\NR
\NC barstyle \NC \NC 下划线样式 \NC\NR
\NC barcolor \NC \NC 下划线颜色 \NC\NR
\NC frame \NC \NC 是否显示边框 \NC\NR
\NC stopper \NC \NC 题号后缀 \NC\NR
\NC distance \NC \NC 选项间距 \NC\NR
\NC before \NC \NC 环境前内容 \NC\NR
\NC after \NC \NC 环境后内容 \NC\NR
\HL
\stoptabulate

\section{进阶用法}

\subsection{closechoice 命令}

\type{\closechoice} 命令用于在完形填空中插入选项，语法为：
\type{\closechoice[选项1,选项2,选项3,...]}。

正确答案使用 \type{[*]} 标记：\type{\closechoice[选项1,{[*]正确选项},选项3]}。

\subsection{optclose 环境}

\type{optclose} 环境是 \type{close} 的变体，用于选项式填空：

\startbuffer
\startoptclose[point=2,showanswer=true]
The latest \ocitem[answer=engineering]{engineer} techniques are applied to create
this protective \ocitem[answer=functional]{function} structure.
\stopoptclose
\stopbuffer

\typebuffer \getbuffer

\chapter{写作区域}

\section{基本用法}

\subsection{作文格}

使用 \type{\writingbox} 创建作文格：

\startbuffer
\writingbox
  [column=6,
   row=6,
   evencolor=blue,
   oddcolor=darkgray,
   evenheight=1cm,
   oddheight=.5cm,
   height=10cm,
   width=\textwidth]{}
\stopbuffer

\typebuffer \getbuffer

\subsection{方格纸}

使用 \type{\writinggrid} 创建方格纸：

\startbuffer
\writinggrid
  [cellwidth=1cm,
   column=12,
   row=8,
   cellcolor=darkgray,
   height=8cm,
   width=\textwidth]{}
\stopbuffer

\typebuffer \getbuffer

\section{参数说明}

\subsection{writingbox 参数}

\showsetup[setupwritingbox]

\starttabulate[|l|l|p|]
\HL
\NC 参数 \NC 默认值 \NC 说明 \NC\NR
\HL
\NC column \NC \NC 列数 \NC\NR
\NC row \NC \NC 行数 \NC\NR
\NC evencolor \NC \NC 偶数行颜色 \NC\NR
\NC oddcolor \NC \NC 奇数行颜色 \NC\NR
\NC evenheight \NC \NC 偶数行高度 \NC\NR
\NC oddheight \NC \NC 奇数行高度 \NC\NR
\NC height \NC \NC 总高度 \NC\NR
\NC width \NC \NC 总宽度 \NC\NR
\NC align \NC \NC 对齐方式 \NC\NR
\NC rulethickness \NC \NC 线条粗细 \NC\NR
\HL
\stoptabulate

\subsection{writinggrid 参数}

\showsetup[setupwritinggrid]

\starttabulate[|l|l|p|]
\HL
\NC 参数 \NC 默认值 \NC 说明 \NC\NR
\HL
\NC column \NC \NC 列数 \NC\NR
\NC row \NC \NC 行数 \NC\NR
\NC cellwidth \NC \NC 单元格宽度 \NC\NR
\NC cellcolor \NC \NC 单元格颜色 \NC\NR
\NC height \NC \NC 总高度 \NC\NR
\NC width \NC \NC 总宽度 \NC\NR
\HL
\stoptabulate

\section{进阶用法}

\subsection{writing 环境}

\type{writing} 环境用于创建作文题目，支持设置分数和答案：

\startbuffer
\startwriting[point=45,showpoint=true]
  作文题目说明
  
  请根据以下要求写作：
  \startitemize
  \item 要求一
  \item 要求二
  \stopitemize
  
  \startanswer
  范文内容...
  \stopanswer
\stopwriting
\stopbuffer

\typebuffer \getbuffer

\subsection{subwriting 环境}

\type{subwriting} 环境用于创建子作文题目：

\startbuffer
\startwriting[point=100]
  作文题目说明
  
  \startsubwriting[point=50]
    子作文题目一
    \startanswer
    范文一...
    \stopanswer
  \stopsubwriting
  
  \startsubwriting[point=50]
    子作文题目二
    \startanswer
    范文二...
    \stopanswer
  \stopsubwriting
\stopwriting
\stopbuffer

\typebuffer \getbuffer

\subsection{writing 参数}

\type{writing} 环境继承自 \type{question} 环境，参数通过 \type{\setupwriting} 设置：

\starttabulate[|l|l|p|]
\HL
\NC 参数 \NC 默认值 \NC 说明 \NC\NR
\HL
\NC point \NC \NC 作文分数 \NC\NR
\NC answer \NC 参见参考作文 \NC 参考答案 \NC\NR
\NC totalanswer \NC \NC 总参考答案 \NC\NR
\NC showanswer \NC \NC 是否显示答案 \NC\NR
\NC showpoint \NC \NC 是否显示分数 \NC\NR
\NC style \NC \NC 题目字体样式 \NC\NR
\NC color \NC \NC 题目颜色 \NC\NR
\NC option \NC n,packed \NC 列表选项 \NC\NR
\NC before \NC \NC 环境前内容 \NC\NR
\NC after \NC \NC 环境后内容 \NC\NR
\HL
\stoptabulate

\chapter{符号定义}

模块内置了多种标记符号（需要安装 NotoSansSymbols2 字体）：

\starttabulate[|l|r|l|r|]
\HL
\NC 命令 \NC 符号 \NC 命令 \NC 符号 \NC\NR
\HL
\NC \type{\symbol{marksquare}} \NC \symbol{marksquare}
\NC \type{\symbol{markcircle}} \NC \symbol{markcircle} \NC\NR
\NC \type{\symbol{markcheck}} \NC \symbol{markcheck}
\NC \type{\symbol{markcross}} \NC \symbol{markcross} \NC\NR
\NC \type{\symbol{markheavycheck}} \NC \symbol{markheavycheck}
\NC \type{\symbol{markheavycross}} \NC \symbol{markheavycross} \NC\NR
\NC \type{\symbol{marksquarecheck}} \NC \symbol{marksquarecheck}
\NC \type{\symbol{marksquarecross}} \NC \symbol{marksquarecross} \NC\NR
\NC \type{\symbol{markcirclecheck}} \NC \symbol{markcirclecheck}
\NC \type{\symbol{markcirclecross}} \NC \symbol{markcirclecross} \NC\NR
\HL
\stoptabulate

\section{圆圈文本}

\type{\circledtext} 命令用于创建带圆圈或方框的文本：

\startbuffer
\circledtext{1}\quad
\circledtext[radius=0.5]{A}\quad
\circledtext[radius=0,fillcolor=black,color=white]{B}\quad
\circledtext[radius=0.2]{C}
\stopbuffer

\typebuffer \getbuffer

\subsection{circledtext 参数}

\showsetup[setupcircledtext]

\starttabulate[|l|l|p|]
\HL
\NC 参数 \NC 默认值 \NC 说明 \NC\NR
\HL
\NC frametype \NC circle \NC 边框类型：circle、square \NC\NR
\NC framecolor \NC \NC 边框颜色 \NC\NR
\NC fillcolor \NC \NC 填充颜色 \NC\NR
\NC dashcolor \NC \NC 虚线颜色 \NC\NR
\NC dashlinewidth \NC \NC 虚线宽度 \NC\NR
\NC color \NC \NC 文本颜色 \NC\NR
\NC radius \NC \NC 圆角半径 \NC\NR
\NC effect \NC normal \NC 特效 \NC\NR
\NC scaled \NC \NC 缩放比例 \NC\NR
\HL
\stoptabulate

\subsection{圆圈编号转换}

模块提供多种圆圈编号转换格式，可用于列表编号：

\startbuffer
\startitemize[Co:num,packed]
  \item 第一项
  \item 第二项
  \item 第三项
\stopitemize

\startitemize[CO:num,packed]
  \item 第一项
  \item 第二项
  \item 第三项
\stopitemize
\stopbuffer

\typebuffer \getbuffer

\starttabulate[|l|l|]
\HL
\NC 转换格式 \NC 说明 \NC\NR
\HL
\NC Co:num \NC 白底黑字圆圈数字 \NC\NR
\NC CO:num \NC 黑底白字圆圈数字 \NC\NR
\NC Cr:num \NC 白底黑字圆角方框数字 \NC\NR
\NC CR:num \NC 黑底白字圆角方框数字 \NC\NR
\NC Cs:num \NC 白底黑字方框数字 \NC\NR
\NC CS:num \NC 黑底白字方框数字 \NC\NR
\NC Co:Alph \NC 大写圆圈字母 \NC\NR
\NC Co:alph \NC 小写圆圈字母 \NC\NR
\NC Co:cn \NC 圆圈中文数字 \NC\NR
\NC Co:hira \NC 圆圈平假名 \NC\NR
\NC Co:kata \NC 圆圈片假名 \NC\NR
\HL
\stoptabulate

\chapter{词典条目}

\section{基本用法}

\type{lexicalentry} 环境用于创建词典风格的条目，适合制作词汇表、术语表或语言学习材料。
每个条目可以包含词性标签、重音标记、释义、例句等信息。

\startbuffer
\startlexicalentry[class=名,accent=1]{太陽}
  \startlexicallists
    \item 太陽系の中心をなし、地球にもっとも近い恒星。お日さま。\antonym{太陰}。
    \item 光明・希望の象徴。「心の\alias」
  \stoplexicallists
\stoplexicalentry
\stopbuffer

\typebuffer \getbuffer

\section{参数说明}

\showsetup[setuplexicalentry]

\starttabulate[|l|l|p|]
\HL
\NC 参数 \NC 默认值 \NC 说明 \NC\NR
\HL
\NC class \NC \NC 词性标签（如"名"、"动"） \NC\NR
\NC accent \NC \NC 重音标记（数字） \NC\NR
\NC number \NC \NC 编号 \NC\NR
\NC style \NC \NC 词条样式 \NC\NR
\NC color \NC \NC 词条颜色 \NC\NR
\NC before \NC \NC 条目前内容 \NC\NR
\NC after \NC \NC 条目后内容 \NC\NR
\HL
\stoptabulate

\section{进阶用法}

\subsection{词典列表}

\type{lexicallists} 是配套的列表环境，支持四级嵌套，每级使用不同的编号样式：

\startbuffer
\startlexicalentry[class=動]{走る}
  \startlexicallists
    \item 足を動かして進む。\antonym{止まる}。
    \item 逃げる。
      \startlexicallists
        \item 時間が経過する。
        \item 機械が動く。
      \stoplexicallists
  \stoplexicallists
\stoplexicalentry
\stopbuffer

\typebuffer \getbuffer

\subsection{词典枚举}

模块提供多种词典枚举环境，用于结构化展示词条信息：

\starttabulate[|l|l|]
\HL
\NC 环境 \NC 说明 \NC\NR
\HL
\NC lexicalgrammar \NC 语法说明 \NC\NR
\NC lexicalmeaning \NC 释义说明 \NC\NR
\NC lexicaljunction \NC 接续说明 \NC\NR
\NC lexicalexample \NC 例句说明 \NC\NR
\NC lexicalcomparing \NC 对比说明 \NC\NR
\NC lexicalsummary \NC 总结说明 \NC\NR
\HL
\stoptabulate

\startbuffer
\startlexicalentry[class=形]{美しい}
  \startlexicalmeaning
    見た目や心がすばらしい。きれいだ。
  \stoplexicalmeaning
  
  \startlexicalexample
    美しい景色を眺める。\\
    美しい心を持つ。
  \stoplexicalexample
  
  \startlexicalcomparing
    「美しい」は芸術的な美しさ、「きれい」は清潔さを強調。
  \stoplexicalcomparing
\stoplexicalentry
\stopbuffer

\typebuffer \getbuffer

\subsection{关联命令}

词典条目支持多种关联命令，用于标记词条间的关系：

\starttabulate[|l|p|]
\HL
\NC 命令 \NC 说明 \NC\NR
\HL
\NC \type{\antonym{词}} \NC 反义词 \NC\NR
\NC \type{\synonym{词}} \NC 同义词 \NC\NR
\NC \type{\alias} \NC 词条别名引用 \NC\NR
\NC \type{\refer{词}} \NC 参见引用 \NC\NR
\NC \type{\example{例}} \NC 例句标记 \NC\NR
\NC \type{\junction{接}} \NC 接续标记 \NC\NR
\HL
\stoptabulate

\chapter{高级功能}

\section{无限层级支持}

模块支持无限层级的题目嵌套，自动管理答案树：

\startbuffer
\startquestion[point=20]
第一层大题
  \startquestion[point=10]
  第二层子题
    \startquestion[point=5]
    第三层小题
    \stopquestion
  \stopquestion
\stopquestion
\stopbuffer

\typebuffer \getbuffer

\section{数据持久化}

答案数据自动保存到 .tuc 文件，可在第二次编译时访问：

\starttyping
% 访问保存的答案数据
\datasetvariable{UnlimitedDataset}{题目ID}{answer}
\datasetvariable{UnlimitedDataset}{题目ID}{point}
\stoptyping

\section{答案树结构}

basicexam 模块使用答案树结构来管理所有题目的答案和分数信息。
每个题目节点包含以下属性：

\starttabulate[|l|p|]
\HL
\NC 属性 \NC 说明 \NC\NR
\HL
\NC id \NC 节点唯一标识（如 \type{jobname-1-Q-1}） \NC\NR
\NC order \NC 层级顺序（如 \type{1}, \type{2-1}, \type{2-2-1}） \NC\NR
\NC answer \NC 答案内容 \NC\NR
\NC point \NC 当前节点分数（叶子节点有效） \NC\NR
\NC totalpoint \NC 该节点及子节点总分 \NC\NR
\NC totalnumber \NC 该节点及子节点总数 \NC\NR
\NC is_leaf \NC 是否为叶子节点 \NC\NR
\NC depth \NC 节点深度（0=根节点） \NC\NR
\HL
\stoptabulate

\section{节点信息获取}

\subsection{当前节点信息}

\starttabulate[|l|p|]
\HL
\NC 命令 \NC 说明 \NC\NR
\HL
\NC \type{\UnlimitedCurrentID} \NC 获取当前节点 ID \NC\NR
\NC \type{\UnlimitedGetDepth} \NC 获取当前节点深度 \NC\NR
\NC \type{\UnlimitedCurrentAnswer} \NC 获取当前节点答案 \NC\NR
\NC \type{\UnlimitedCurrentPoint} \NC 获取当前节点总分 \NC\NR
\HL
\stoptabulate

\subsection{指定节点信息}

使用 \type{\UnlimitedGetAnswer} 获取指定节点的属性：

\startbuffer
% 获取题目1的分数
\UnlimitedGetAnswer{jobname-1-Q-1}{point}

% 获取题目1的答案
\UnlimitedGetAnswer{jobname-1-Q-1}{answer}

% 获取题目1的总分（包含子题）
\UnlimitedGetAnswer{jobname-1-Q-1}{totalpoint}

% 获取题目1的总数（包含子题）
\UnlimitedGetAnswer{jobname-1-Q-1}{totalnumber}
\stopbuffer

\typebuffer

\section{调试功能}

在 check 模式下，模块会自动输出调试信息：

\starttyping
\usemodule[basicexam][mode={teacher,check}]
\stoptyping

也可以手动调用调试命令：

\starttyping
% 输出答案树结构
\UnlimitedDebug

% 导出答案到 TSV 文件
\UnlimitedExport{answers.tsv}
\stoptyping

\chapter{贡献与修改}

\section{提交需求}

如果您在使用过程中发现 bug 或有新功能需求，可以通过以下方式提交：

\startitemize[packed]
\item 在 GitHub 仓库提交 Issue
\item 发送邮件至作者邮箱
\item 在相关论坛发帖讨论
\stopitemize

提交需求时，请尽量提供以下信息：

\startitemize[packed]
\item ConTeXt 版本信息（运行 \type{context --version}）
\item 最小复现示例
\item 期望的行为和实际的行为
\stopitemize

\section{修改本模块}

本模块采用 LuaMetaTeX 编写，源代码位于 \type{t-basicexam.mklx} 文件中。

\subsection{模块结构}

\starttabulate[|l|p|]
\HL
\NC 文件 \NC 说明 \NC\NR
\HL
\NC t-basicexam.mklx \NC 主模块文件，包含所有命令定义 \NC\NR
\NC i-basicexam.xml \NC 接口定义文件，用于命令文档生成 \NC\NR
\NC doc/basicexam-manual.tex \NC 使用手册 \NC\NR
\NC doc/test-exam.tex \NC 测试文件 \NC\NR
\HL
\stoptabulate

\subsection{修改建议}

\startitemize[packed]
\item 修改前请备份原文件
\item 建议先阅读源代码注释
\item 修改后请运行测试文件验证
\item 如有重大改进，欢迎提交 Pull Request
\stopitemize

\stoptext
